\subsection{I-41. Weight-two connectors and the balanced ten-variable theorem}
\label{I-41}

\paragraph{Statement and scope.}
Let $E$ be any field, and fix arbitrary distinct nodes
$a_1,\ldots,a_n$ in $E$. Every nonzero homogeneous multilinear
polynomial $f$ of width at most three in some variable order,
satisfying
\[
 f(X+s\mathbf1)=f(X),\qquad 0\le n-2\deg f\le2,
 \tag{I.131}
\]
has nonzero pullback under the balanced rational map
\[
 R_i=t+\frac{z_1}{a_i-y_1}+\frac{z_2}{a_i-y_2}.
 \tag{I.132}
\]
The order is arbitrary and unknown to the map, and the result
holds in every characteristic. Together with P-17 and I-40,
this proves the same balanced two-pair conclusion for every
nonzero multilinear width-at-most-three polynomial on at most
ten variables, without a homogeneity or translation-invariance
promise.

For total degree at most $D$, both assertions admit the finite
sample bound
\[
 (n+1)^2\binom{D+3}{3}
 \tag{I.133}
\]
over a working field of size at least
$\max\{2n+1,D+1\}$, with the extension-input convention of I-40.
The new weight-two argument uses the common rational center
$t\mathbf1$. It does not claim arbitrary nonzero center slopes
or the ordinary unweighted centered-$\mathrm{SV}_2$ conclusion
for this entire larger class.

\paragraph{1. Reduction to an irreducible tensor of weight two.}
Work over an algebraic closure as a proof device. Suppose a
nonzero input in (I.131) vanishes under (I.132). Since the
substitution target is a field, some nonconstant irreducible
factor $g$ also vanishes. Irreducible factors of a multilinear
polynomial use disjoint variable sets; specializing the other
sets to nonzero scalar factor values realizes $g$, up to scalar,
as a width-at-most-three polynomial in the induced order, as
proved in I-37.

These factors are homogeneous. Indeed the minimum and maximum
total degrees in a product add, and equality of those degrees
for a homogeneous product forces equality for every factor.
Each factor is also invariant under common translation of its
own variables. To see this, divide the translated product
identity by the original product over $E(s)$. The resulting
factors are rational functions in disjoint sets of variables,
so each ratio belongs to $E(s)^\times$. Comparing its highest
homogeneous part with the original homogeneous factor forces
that ratio to be one. I-40 therefore gives a nonnegative weight
for every factor. Their weights sum to $n_{\rm used}-2\deg f$,
which is at most the original $n-2\deg f$; unused coordinates
only increase the latter. Thus the chosen factor has weight
at most two.

Weights zero and one are already covered by I-40, as are all
inputs on at most nine variables. We may henceforth assume
that $f$ is irreducible, depends on every remaining variable,
has weight $m=n-2D=2$, and has $n\ge10$. Construct its actual
minimal affine cores by I-40's support inclusions, and pad
consistently to width three. We use those new minimal cores,
not a possibly nonminimal representation inherited from the
original factorization. Write the unpadded bond ranks as $r_i$.
Every interior cut has $r_i\ge2$: rank one would split the
scalar into a product in two nonempty disjoint variable sets.
Moreover
\[
 r_1=r_{n-1}=2,
 \tag{I.134}
\]
since a single multilinear variable supplies at most two
coefficient vectors. In particular the first and last internal
factors, $A_2$ and $A_{n-1}$, are singular after padding.

As in I-37, every strict internal interval has generic rank at
least two. A rank-one interval has a polynomial outer
factorization over the multivariate UFD and splits the scalar
into two nonconstant factors, because the first and last
variables occur on opposite sides. Each such interval therefore
has rank two or three. I-40 bounds every run of rank-three
factors by two. The only new rank-preservation issue is a pair
of consecutive rank-three cores between rank-two anchors.

\paragraph{2. A saturated run has invariant bounding support spaces.}
Suppose $A_{j+1},A_{j+2}$ are both rank three. They are a
maximal such run. The determinant characters of I-40 satisfy
\[
 \delta_j=0,\qquad \delta_{j+1}=3,
 \qquad\delta_{j+2}=6,
 \qquad\delta'_{j+2}=0.
 \tag{I.135}
\]
Indeed both increments are three and the allowed interval is
$0\le\delta_i\le3m=6$ at these cuts.

A highest-weight line of weight zero is fixed by the whole
$\mathrm{SL}_2$. For completeness, the Bruhat identity in the
I-40 audit shows that its opposite-unipotent orbit is a
polynomial of degree at most zero; it is therefore constant.
The two unipotent subgroups generate the group. Apply this to
the determinant lines in (I.135). A subspace is determined by
its nonzero decomposable exterior vector, equivalently by the
vectors whose exterior product with that vector is zero.
Thus $L_j$ and $R_{j+2}$ themselves are $\mathrm{SL}_2$-stable
three-dimensional spaces.

Taking all single-coordinate contractions propagates this
stability to $L_{j-1}$ and to $R_{j+3}$, respectively. These
are the actual support spaces: contracting $L_j$ by all
functionals on its final physical factor spans $L_{j-1}$,
and the analogous statement holds on the right. Every torus
weight in $L_j\subseteq V^{\otimes j}$ is congruent to $j$
modulo two. Since its determinant weight is zero and its
dimension is three, $j$ is even. Also $n$ is even because
$n-2D=2$. The adjacent support spaces lie in odd tensor powers.
Weyl pairing of their nonzero integer weights makes their
dimensions even; the width bound and nonzeroness therefore give
\[
 \dim L_{j-1}=\dim R_{j+3}=2.
 \tag{I.136}
\]
This also shows that both adjacent anchors exist away from the
scalar endpoints.

\paragraph{3. The canonical anchors, without a semisimplicity assumption.}
Let $W=L_{j-1}$. Its two torus weights are $q,-q$ for an odd
positive integer $q$. If $q>1$, all four weights of $W\otimes V$
are nonzero, namely $q+1,q-1,-q+1,-q-1$. Every
$\mathrm{SL}_2$-stable subspace of this tensor product then has
even dimension, again by Weyl pairing. This contradicts its
three-dimensional stable subspace $L_j$. Therefore $q=1$.
This elementary argument excludes nontrivial Frobenius twists
as well as every other possible larger character.

A two-dimensional rational $\mathrm{SL}_2$-module with weights
$1,-1$ is the standard module in every characteristic. One can
check the convention directly: torus conjugation forces the
upper and lower unipotent entries to be $b s$ and $c s$;
neither is zero, since the two unipotents generate the group.
The Gauss identity for their product forces $bc=1$, and a
basis rescaling gives the standard matrices. This is also the
degree-one case of the established I-B09 input.

Now identify $W$ with $V$. A stable three-plane in $V\otimes V$
must contain the one-dimensional weight-two space: otherwise
Weyl symmetry would remove weight minus two as well, leaving
only the two-dimensional weight-zero space. Applying the
opposite unipotent to $e_0\otimes e_0$ shows that the three-plane
contains, and therefore equals,
\[
 \Gamma^2 V=\operatorname{span}\{
 e_0e_0,\ e_0e_1+e_1e_0,\ e_1e_1\}.
 \tag{I.137}
\]
Here juxtaposition denotes tensor product. The three displayed
vectors are independent in characteristic two as well. We do
not identify this submodule with its dual or invoke a
semisimple decomposition of $V\otimes V$.

Using its displayed weight basis, the left anchor is exactly
\[
 C(u)=\begin{pmatrix}1&u&0\\0&1&u\end{pmatrix},
 \qquad u=X_j.
 \tag{I.138}
\]
The same argument on the right identifies $R_{j+2}$ with the
three-plane (I.137) in $V\otimes R_{j+3}$, where the latter
two-dimensional residual space is standard.

The two full-rank cores have nonzero constant determinants by
I-40. Make their constant matrices the identity by constant
changes of the two intermediate bases, starting from the chosen
left weight basis. The torus actions on those bases are shifted
successively by weights one and two. Their linear coefficient
matrices consequently have entries only one step above the
diagonal: their row weight minus column weight must equal two.
Thus their product has the form
\[
 M(x,y)=
 \begin{pmatrix}1&A&B\\0&1&C\\0&0&1\end{pmatrix},
 \tag{I.139}
\]
where $A,C$ are homogeneous linear forms in $x=X_{j+1}$ and
$y=X_{j+2}$, and $B$ is homogeneous quadratic (in fact a scalar
multiple of $xy$). All these changes of basis are over the
coefficient field and are independent of the evaluation-center
parameters.

The left weights at the end of the run are $(4,2,0)$, while
the right weights in $R_{j+2}$ are $(2,0,-2)$. The tensor has
weight two and full rank three across that cut. Its coefficient
pairing is therefore antidiagonal with three nonzero coefficients
$d_1,d_2,d_3$. Expand the canonical right three-plane (I.137),
and order the next left residual basis oppositely to the right
standard basis. The right anchor is
\[
 D(v)=\begin{pmatrix}d_1v&0\\d_2&d_2v\\0&d_3\end{pmatrix},
 \qquad v=X_{j+3},\qquad d_1d_2d_3\ne0.
 \tag{I.140}
\]
This uses the actual dual pairing. In particular it remains
valid when $\Gamma^2V$ is not self-dual in characteristic two.
The full-rank constant gauges identify torus actions only;
they do not assert an isomorphism of the neighboring
upper-unipotent modules.

\paragraph{4. The four-coordinate connector cannot vanish at one pole.}
The connecting matrix between the two-dimensional cuts is
$C(u)M(x,y)D(v)$. Its determinant is
\[
 \begin{split}
 \det(CMD)={}&d_1d_2v^2+d_1d_3v(C+u)\\
 &+d_2d_3\bigl(u^2+Au+AC-B\bigr).
 \end{split}
 \tag{I.141}
\]
This follows by multiplying the two-by-three and three-by-two
matrices; equivalently its $u^2$ coefficient is the bottom
two-row minor of $MD$. It is a homogeneous quadratic in the
four distinct coordinate variables, and its $u^2$ coefficient
is the constant $d_2d_3\ne0$.

Set $\alpha_i=t+z/(a_i-y_0)$, where $t,z,y_0$ are independent.
At $y_0=a_j$, only $u=\alpha_j$ has a pole among these four
coordinates. The term $d_2d_3u^2$ has a double pole with nonzero
leading coefficient $d_2d_3z^2$, and every other term in
(I.141) has pole order at most one there. Hence the connector
determinant survives this substitution in every characteristic.

We now verify all strict internal interval ranks at this same
one-pole center. Every individual full-rank core remains
invertible; its determinant is a nonzero constant. Every
rank-two core keeps rank two, because some nonzero univariate
minor is evaluated at a transcendental coordinate. In the
second compound word, these are respectively full-rank and
rank-one factors. A rank-two interval has a nonzero compound
product. Factoring each rank-one compound by a nonzero
univariate pivot writes it as nonzero endpoint vectors times
the scalar connectors between consecutive rank-one factors.

If no or one full-rank factor lies between two such factors,
the connector depends on at most three coordinates and survives
by I-38's three-coordinate dominance. If two lie between them,
they form exactly the saturated run just analyzed. At its
adjacent two-dimensional cuts the compound connector is the
determinant (I.141), up to nonzero univariate pivot factors,
so it too survives. Thus every compound interval product that
was nonzero stays nonzero. A rank-three interval consists
entirely of invertible factors and stays invertible. All
strict internal interval ranks are therefore preserved
simultaneously. This asserts preservation of ranks, not of
every individual minor.

\paragraph{5. Endpoints and compatible constant gauges.}
Absorb the first and last factors into the affine row and column.
By (I.134), the first and last internal factors are rank two.
The polynomial row $A_1(X_1)A_2(X_2)$ is nonzero, since otherwise
the full scalar would be zero. Pairwise dominance preserves a
nonzero entry at the center. The same holds for
$A_{n-1}(X_{n-1})A_n(X_n)$. Propagation through the remaining
rank-preserved constant intervals is injective on the persistent
two-dimensional active spaces, so both endpoint active
projections are nonzero. No four-coordinate endpoint assertion
is being assumed.

The full constant interior product has rank two, and every
subinterval has rank two or three. The compatible-space proof
of I-38, part~4, applies to this rank profile without any bound
on run length: propagate a two-dimensional complement of the
full-product row kernel, and take the suffix preimages of a
final complementary line. Each passive line is either the
rank-two suffix kernel or the inverse image under an invertible
suffix. In the transported bases the active blocks are
$I_2$, the singular passive blocks are zero, and the invertible
passive blocks can be scaled to one.

The endpoint active vectors pair to $f(\boldsymbol\alpha)=0$.
Normalize them to $e_1^{\mathsf T}$ and $e_2$, and remove the
passive endpoint components by the upper and lower shears of
I-37. They leave the first and last singular constants fixed;
if those coincided the same identities would permit both
shears. Over $K=\overline E(t,z,y_0)$, with independent residual
variables $Y_i=X_i-\alpha_i$, the resulting scalar is
\[
 f(\boldsymbol\alpha+\mathbf Y)=
 (e_1^{\mathsf T}+Y_1p)
 \prod_{i=2}^{n-1}(C_i+Y_iN_i)
 (e_2+Y_nq),
 \quad C_i\in\{P,I_3\},\quad P=\operatorname{diag}(1,1,0).
 \tag{I.142}
\]
The $I_3$ runs have length at most two, and $C_2=C_{n-1}=P$.
Pole extraction from the assumed two-pole identity, exactly
as in I-37, gives $f_i(\boldsymbol\alpha)=0$ for every $i$.
It follows that
\[
 p_2=q_1=(N_i)_{12}=0.
 \tag{I.143}
\]
At each singular factor the coefficient of $Y_i$ in
$\det(P+Y_iN_i)=0$ also gives $(N_i)_{33}=0$.
No such assertion is made at invertible factors.

\paragraph{6. Three tangent directions kill the bandwidth-three Hessian.}
Let $H=(f_{ij}(\boldsymbol\alpha))$, and put
$h_i=1/(a_i-y_0)$. Differentiate the already zero vector
$(f_i(\boldsymbol\alpha))$ with respect to $t,z,y_0$.
Equivalently take first coefficients under separate parameter
increments. Since $z\ne0$ in $K$, this gives
\[
 H\mathbf1=H\mathbf h=H\mathbf h^{\circ2}=0.
 \tag{I.144}
\]
All three identities hold in every characteristic. The diagonal
of $H$ is zero by multilinearity.

If a singular position lies strictly between $i$ and $j$, the
intervening constant product contains $P$. It removes the
passive state, while the active derivative blocks in (I.143)
are lower triangular, so $H_{ij}=0$. Otherwise all those
constants are identities, and $H_{ij}$ is exactly the product
of the state-one-to-passive derivative entry at $i$ and the
passive-to-state-two derivative entry at $j$. The endpoint
versions are $p_3$ and $q_3$. These statements follow directly
from (I.142); the other active contributions contain a zero
$(1,2)$ entry. Since no invertible run has length three,
\[
 H_{ij}=0\qquad\text{when }|i-j|>3.
 \tag{I.145}
\]

The first row of $H$ has at most three remaining entries.
For three distinct future indices $j,k,\ell$, the three
equations (I.144) have coefficient determinant
\[
 \det\begin{pmatrix}
 1&1&1\\h_j&h_k&h_\ell\\h_j^2&h_k^2&h_\ell^2
 \end{pmatrix}
 =(h_k-h_j)(h_\ell-h_j)(h_\ell-h_k)\ne0.
 \tag{I.146}
\]
The one- and two-entry cases use the corresponding initial
rows of this Vandermonde matrix. Thus the first row is zero.
By symmetry, the earlier entries of the next row are now zero,
leaving at most three future entries to which the same argument
applies. Induction proves $H=0$.

Finally expand (I.142) in state paths. Direct passage from
active state one to active state two is excluded by (I.143).
Before a path first reaches active state two, take its last
departure from active state one into the passive state; if it
starts passively, take its passive endpoint entry instead.
Any intervening passive stays must occur at invertible factors,
because a singular factor has both its passive constant and
linear entry zero. The product of the departure and arrival
entries is precisely the corresponding $H_{ij}=0$.
Additional passive stays only multiply this zero product.
The same reasoning applies to a final output through $q_3$;
an output through active state one is excluded by $q_1=0$.
Therefore every output path is zero, contradicting the
nonzero polynomial $f$. This proves (I.132) under (I.131).

\paragraph{7. All multilinear inputs on at most ten variables.}
I-40 already covers at most nine variables. For ten variables,
the established balanced-ideal degree interval is $3\le d\le7$.
The cubic theorem of P-17 excludes degree three, and reversal
excludes degree seven. Consider a homogeneous quartic
annihilator of width at most three. Its common translation has
the form $f+s g$ because the ideal has no components of degree
at most two. A nonzero $g$ would be a cubic annihilator with
the same width bound: as in I-40, it is the leading coefficient
of a product of invertible local translations, and the
coefficient-cut four-minors remain zero at the leading limit.
P-17 forces $g=0$. The quartic is translation-invariant of
weight $10-2\cdot4=2$, so the theorem just proved excludes it.
Reversal also excludes degree six.

For a homogeneous quintic annihilator, every positive-power
translation coefficient has degree at most four and remains
in the balanced ideal. If any were nonzero, the coefficient
of the highest occurring power would retain width at most
three by the same limit argument. The already excluded cubic
and quartic cases, and the absence of lower degrees, rule this
out. The quintic is therefore translation-invariant of weight
zero, and I-40 excludes it as well. This covers every allowed
homogeneous degree. Passing to a lowest nonzero homogeneous
component, which preserves the ROABP width in the same order,
proves the general multilinear ten-variable assertion.

The finite balanced sampler is exactly I-40, part~6: multiply
the pullback by $\prod_i(a_i-y_1)(a_i-y_2)$, interpolate its two
selector degrees at most $n$ away from every node, and its
three amplitude parameters of total degree at most $D$ on
the Newton simplex. This proves (I.133) and its field scope
for both classes.\hfill$\square$

\paragraph{Provenance and the next obstruction.}
The proof uses the established representation input I-B09 and
the fully proved I-37--I-40 normalization and rank arguments.
The purity agent independently checked the canonical
$\Gamma^2$ anchors, the actual dual pairing in characteristic
two, the connector determinant and the endpoint simplification.
Those checks do not assume that $\Gamma^2V$ is self-dual or
that constant full-rank gauges intertwine the entire Borel
action. The present result makes no historical priority claim.

At weight three, I-40 permits an invertible run of length three.
The three tangent directions in (I.144) do not by themselves
kill a bandwidth-four symmetric matrix. Shorter nonsaturated
runs also need not have both invariant bounding support spaces.
The general mixed affine two-pair gate, the unrestricted
eleven-variable case, and polynomial dependence on a growing
ROABP width remain separate questions.
